% from ega4-10-fr.tex
\oldpage[IV-3]{108}
est local~; donc (5.6.3, (iv)) $\mathcal O_{X,x}$ est universellement caténaire. Pour
démontrer que $X$ vérifie la condition $3^\circ$, considérons plusieurs cas~:

\emph{a)} $X$ est un sous-préschéma fermé irréductible de $S$~; soient $S'$ une
composante irréductible de $S$ contenant $X$, $\xi$ le point générique de $X$,
$x$ un point fermé de $X$~; pour tout $s\in S'$, $\mathcal O_{S',s}$, quotient de
$\mathcal O_{S,s}$, est caténaire (5.6.1), donc les conditions $2^\circ$ et $3^\circ$
entraînent que $S'$ est biéquidimensionnel ($\mathbf 0$, 14.3.3). En vertu de
(5.1.2) et de ($\mathbf 0$, 14.3.3.2), on a donc
\[
  \dim(\mathcal O_{X,x})=\dim(\mathcal O_{S',x})-\dim(\mathcal O_{S',\xi})
  =\dim(S')-\dim(\mathcal O_{S',\xi})
\]
ce qui prouve que $\dim(\mathcal O_{X,x})$ ne dépend pas du point fermé $x$
considéré, d'où l'assertion dans ce cas (5.1.4).

\emph{b)} $X$ est irréductible et $g$ dominant. Alors, pour tout point fermé $x$
de $X$, $g(x)$ est fermé dans $S$ par (10.4.7)~; comme $\mathcal O_{S,g(x)}$ est
universellement caténaire, il résulte de (5.6.5.3) que l'on a
\[
  \dim(\mathcal O_{X,x})=\dim(\mathcal O_{S,g(x)})+e
\]
où $e=\dim(g^{-1}(\zeta))$, $\zeta$ étant le point générique de $S$. Comme
$\dim(S)=\dim(\mathcal O_{S,g(x)})$ en vertu de la condition $3^\circ$ pour $S$, on
a $\dim(\mathcal O_{X,x})=\dim(S)+e$ pour tout point fermé $x\in X$~; cela prouve
la condition $3^\circ$ pour $X$ (5.1.4), et en même temps la formule
\begin{equation}
\label{IV.10.6.1.4-fr}
\tag{10.6.1.4}
\dim(X)=\dim(S)+e.
\end{equation}

\emph{c)} Cas général. En considérant un sous-préschéma réduit $X'$ de $X$ ayant
pour espace sous-jacent une composante irréductible de $X$, on est ramené au cas
où $X$ est intègre~; utilisant (\textbf{I}, 5.2.2) et le cas \emph{a)} prouvé
ci-dessus, on peut ensuite remplacer $S$ par le sous-préschéma réduit ayant
$\overline{g(X)}$ pour espace sous-jacent~; on est alors ramené au cas \emph{b)},
et cela achève de démontrer (i).

(ii) Le morphisme $f$ étant localement de type fini (\textbf{I}, 1.3.4), on peut
appliquer les résultats de (i) en remplaçant $S$ par $Y$~; en outre, comme $X$
est irréductible et $f$ dominant, on peut aussi remplacer $S$ par $Y$ dans
(10.6.1.4), ce qui donne (10.6.1.1).

(iii) L'assertion relative au cas où $\dim(f^{-1}(y))\leq n$ pour tout $y\in Y$
a déjà été démontrée sous des hypothèses plus générales dans (5.6.7). Supposons
que $\dim(f^{-1}(y))\geq n$ pour tout $y\in Y$, et considérons un point générique
$\eta$ d'une composante irréductible $Y'$ de $Y$~; il existe au moins une
composante irréductible $Z$ de $f^{-1}(\eta)$ de dimension $\geq n$~; si $\xi$
est le point générique de $Z$, $\xi$ est aussi point générique d'une composante
irréductible $X'$ de $X$ telle que $f(X')$ soit dense dans $Y'$ ($\mathbf 0_{\mathrm I}$,
2.1.8). Considérons les sous-préschémas réduits de $X$, $Y$ ayant respectivement
pour espaces sous-jacents $X'$, $Y'$, et la restriction $X'\to Y'$ de $f$
(\textbf{I}, 5.2.2)~; il résulte alors de (ii) que
$\dim(X')\geq\dim(Y')+n$, et a fortiori $\dim(X)\geq\dim(Y')+n$~; ceci étant vrai
pour toute composante irréductible $Y'$ de $Y$, on en conclut l'inégalité
(10.6.1.2).

\begin{corollary}[10.6.2]
\label{IV.10.6.2-fr}
Supposons que $X$ vérifie les conditions $1^\circ$, $2^\circ$ et $3^\circ$ de
(10.6.1). Alors, pour tout ouvert $U$ dense dans $X$, on a
$\dim(U)=\dim(X)$.
\end{corollary}

